% GENERATED FILE. Edit canonical lessons, then run npm run generate:books.
% canonical-source-hash: fnv1a64:a3ced5b928523ecb
% canonical-lessons: GE-C287-gestehen, GE-C287-gestaendnis, GE-C287-gefaengnis, GE-C287-verschwinden, GE-C287-fliehen

\chapter{Das Gefängnis}
\label{ch:ge-das-gefangnis}
\hlchaptermodality{287}

\begin{chapteropening}
\textbf{By the end of this chapter:} \emph{I can say how a crime story ends: to confess, a confession, prison, to disappear and to flee.}

Complete the last lesson of chapter 287: I can say how a crime story ends: to confess, a confession, prison, to disappear and to flee.
\end{chapteropening}

\section[gestehen]{gestehen — to confess, to admit}

\label{lesson:GE-C287-gestehen}



\begin{quote}
Before the new one: say the German for guilty, then the German for innocent.
\end{quote}

\subsection*{You'll want to know: gestehen}
\textbf{gestehen} — \textquotedblleft{}to confess, to admit\textquotedblright{}.

\begin{quote}
Nach drei Stunden gestand er alles.
\end{quote}

\begin{quote}
\emph{After three hours he confessed everything.}
\end{quote}

Strong verb: \emph{gestand}, \emph{hat gestanden}.

\begin{cousinweb}
\emph{gestehen} is built on \emph{stehen} \textquotedblleft{}to stand\textquotedblright{}, from Old High German \emph{stān}, Proto-Germanic \emph{stāną} (EtymDB), with \emph{ge-}: to stand by what one did.
\end{cousinweb}

\begin{grammarlens}[title={What you've built}]
One more word for telling, explaining and describing at B1.
\end{grammarlens}

\subsection*{Your turn}
\begin{itemize}
\raggedright
  \item \emph{Say it:} \emph{gestehen}
  \item \emph{Say it:} \emph{gestehen}, once more
  \item \emph{Say it:} a sentence of your own with \emph{gestehen}
  \item \emph{Recall:} say the German for guilty, then the German for innocent, then say \emph{gestehen} again
\end{itemize}

\begin{tcolorbox}[breakable,colback=teal!4,colframe=teal!35!black,title={Before you move on}]
What does gestehen mean? (\textquotedblleft{}To confess, to admit\textquotedblright{}.) Say it once more.
\end{tcolorbox}
\section[das Geständnis]{das Geständnis — the confession}

\label{lesson:GE-C287-gestaendnis}



\begin{quote}
Before the new one: say the German for innocent, then the German for to confess.
\end{quote}

\subsection*{You'll want to know: das Geständnis}
\textbf{das Geständnis} — \textquotedblleft{}the confession\textquotedblright{}.

\begin{quote}
Sein Geständnis überraschte alle.
\end{quote}

\begin{quote}
\emph{His confession surprised everyone.}
\end{quote}

Plural: \emph{die Geständnisse}.

\begin{cousinweb}
\emph{Geständnis} is built on \emph{gestehen} \textquotedblleft{}to confess\textquotedblright{}, from \emph{stehen} \textquotedblleft{}to stand\textquotedblright{} (EtymDB), with the noun ending \emph{-nis}.
\end{cousinweb}

\begin{grammarlens}[title={What you've built}]
One more word for telling, explaining and describing at B1.
\end{grammarlens}

\subsection*{Your turn}
\begin{itemize}
\raggedright
  \item \emph{Say it:} \emph{das Geständnis}
  \item \emph{Say it:} \emph{das Geständnis}, once more
  \item \emph{Say it:} a sentence of your own with \emph{das Geständnis}
  \item \emph{Recall:} say the German for innocent, then the German for to confess, then say \emph{das Geständnis} again
\end{itemize}

\begin{tcolorbox}[breakable,colback=teal!4,colframe=teal!35!black,title={Before you move on}]
What does das Geständnis mean? (\textquotedblleft{}The confession\textquotedblright{}.) Say it once more.
\end{tcolorbox}
\section[das Gefängnis]{das Gefängnis — the prison}

\label{lesson:GE-C287-gefaengnis}



\begin{quote}
Before the new one: say the German for to confess, then the German for the confession.
\end{quote}

\subsection*{You'll want to know: das Gefängnis}
\textbf{das Gefängnis} — \textquotedblleft{}the prison\textquotedblright{}.

\begin{quote}
Er musste zwei Jahre ins Gefängnis.
\end{quote}

\begin{quote}
\emph{He had to go to prison for two years.}
\end{quote}

Plural: \emph{die Gefängnisse}. \emph{Im Gefängnis sitzen} — to be in prison.

\begin{cousinweb}
\emph{Gefängnis} is built on \emph{fangen} \textquotedblleft{}to catch\textquotedblright{}, from Old High German \emph{fāhan}, Proto-Germanic \emph{fanhaną} (EtymDB): the place where the caught are kept.
\end{cousinweb}

\begin{grammarlens}[title={What you've built}]
One more word for telling, explaining and describing at B1.
\end{grammarlens}

\subsection*{Your turn}
\begin{itemize}
\raggedright
  \item \emph{Say it:} \emph{das Gefängnis}
  \item \emph{Say it:} \emph{das Gefängnis}, once more
  \item \emph{Say it:} a sentence of your own with \emph{das Gefängnis}
  \item \emph{Recall:} say the German for to confess, then the German for the confession, then say \emph{das Gefängnis} again
\end{itemize}

\begin{tcolorbox}[breakable,colback=teal!4,colframe=teal!35!black,title={Before you move on}]
What does das Gefängnis mean? (\textquotedblleft{}The prison\textquotedblright{}.) Say it once more.
\end{tcolorbox}
\section[verschwinden]{verschwinden — to disappear, to vanish}

\label{lesson:GE-C287-verschwinden}



\begin{quote}
Before the new one: say the German for the confession, then the German for the prison.
\end{quote}

\subsection*{You'll want to know: verschwinden}
\textbf{verschwinden} — \textquotedblleft{}to disappear, to vanish\textquotedblright{}.

\begin{quote}
Der Dieb verschwand im Wald.
\end{quote}

\begin{quote}
\emph{The thief disappeared into the forest.}
\end{quote}

Strong, with \emph{sein}: \emph{verschwand}, \emph{ist verschwunden}.

\begin{cousinweb}
\emph{verschwinden} is inherited from Middle High German \emph{verswinden}, Old High German \emph{farswindan} (EtymDB), built on an old verb \emph{swindan} \textquotedblleft{}to fade\textquotedblright{}.
\end{cousinweb}

\begin{grammarlens}[title={What you've built}]
One more word for telling, explaining and describing at B1.
\end{grammarlens}

\subsection*{Your turn}
\begin{itemize}
\raggedright
  \item \emph{Say it:} \emph{verschwinden}
  \item \emph{Say it:} \emph{verschwinden}, once more
  \item \emph{Say it:} a sentence of your own with \emph{verschwinden}
  \item \emph{Recall:} say the German for the confession, then the German for the prison, then say \emph{verschwinden} again
\end{itemize}

\begin{tcolorbox}[breakable,colback=teal!4,colframe=teal!35!black,title={Before you move on}]
What does verschwinden mean? (\textquotedblleft{}To disappear, to vanish\textquotedblright{}.) Say it once more.
\end{tcolorbox}
\section[fliehen]{fliehen — to flee, to escape}

\label{lesson:GE-C287-fliehen}



\begin{quote}
Before the new one: say the German for the prison, then the German for to disappear.
\end{quote}

\subsection*{You'll want to know: fliehen}
\textbf{fliehen} — \textquotedblleft{}to flee, to escape\textquotedblright{}.

\begin{quote}
Der Täter floh mit dem Auto.
\end{quote}

\begin{quote}
\emph{The offender fled by car.}
\end{quote}

Strong, with \emph{sein}: \emph{floh}, \emph{ist geflohen}. \emph{Vor etwas fliehen} — to flee from something.

\begin{cousinweb}
\emph{fliehen} is an inherited Germanic verb, usually traced to Proto-Germanic \emph{fleuhaną}; EtymDB lists English \emph{flee} and Dutch \emph{vlieden} as its cousins.
\end{cousinweb}

\begin{grammarlens}[title={What you've built}]
One more word for telling, explaining and describing at B1.
\end{grammarlens}

\subsection*{Your turn}
\begin{itemize}
\raggedright
  \item \emph{Say it:} \emph{fliehen}
  \item \emph{Say it:} \emph{fliehen}, once more
  \item \emph{Say it:} a sentence of your own with \emph{fliehen}
  \item \emph{Recall:} say the German for the prison, then the German for to disappear, then say \emph{fliehen} again
\end{itemize}

\begin{tcolorbox}[breakable,colback=teal!4,colframe=teal!35!black,title={Before you move on}]
What does fliehen mean? (\textquotedblleft{}To flee, to escape\textquotedblright{}.) Say it once more.
\end{tcolorbox}
